\documentclass[11pt]{article}
\usepackage{amsmath,amssymb,amsthm}
\usepackage{lmodern}
\usepackage{geometry}
\usepackage{hyperref}
\usepackage[parfill]{parskip}

\geometry{margin=1.2in}

\newtheorem{proposition}{Proposition}
\theoremstyle{remark}
\newtheorem{remark}{Remark}

\title{%
  \textbf{Partialling Out and the Propensity Score}\\[6pt]
  \large Why $\tilde{D}_i = D_i - p(X_i)$\\[4pt]
  \normalsize ECON 712 --- Supplementary Notes
}
\author{}
\date{}

\begin{document}
\maketitle

%--------------------------------------------------------------------
\section{Setup}
%--------------------------------------------------------------------

Let $D_i \in \{0,1\}$ be a binary treatment and $X_i \in \{0,1\}$ a binary
covariate. Define the propensity score $p(x) = E[D_i \mid X_i = x]$.

We want to show that partialling out $X_i$ from $D_i$ via OLS yields the
residual $\tilde{D}_i = D_i - p(X_i)$.

%--------------------------------------------------------------------
\section{Step 1: OLS Fitted Values in the Saturated Regression}
%--------------------------------------------------------------------

Regress $D_i$ on a constant and $X_i$:
\[
  D_i = \alpha + \gamma X_i + \text{residual}.
\]
Because $X_i$ is binary this regression is \emph{saturated} --- it has one
free parameter per cell.  The OLS estimates are exactly the cell means:
\begin{align*}
  \hat\alpha &= E[D_i \mid X_i = 0] = p(0), \\
  \hat\gamma &= E[D_i \mid X_i = 1] - E[D_i \mid X_i = 0] = p(1) - p(0).
\end{align*}
The fitted value for observation $i$ is therefore
\[
  \hat{D}_i = p(0) + \bigl(p(1)-p(0)\bigr)X_i.
\]

%--------------------------------------------------------------------
\section{Step 2: The Fitted Value Equals $p(X_i)$}
%--------------------------------------------------------------------

Check each cell:
\begin{align*}
  X_i = 0: \quad & \hat{D}_i = p(0) + \bigl(p(1)-p(0)\bigr)\cdot 0 = p(0) = p(X_i). \\
  X_i = 1: \quad & \hat{D}_i = p(0) + \bigl(p(1)-p(0)\bigr)\cdot 1 = p(1) = p(X_i).
\end{align*}
In both cells $\hat{D}_i = p(X_i)$, so the two-coefficient saturated regression
perfectly recovers the propensity score.

%--------------------------------------------------------------------
\section{Step 3: The Residual is $\tilde{D}_i = D_i - p(X_i)$}
%--------------------------------------------------------------------

The OLS residual (the partialled-out treatment) is
\[
  \tilde{D}_i \;=\; D_i - \hat{D}_i \;=\; D_i - p(X_i).
\]

%--------------------------------------------------------------------
\section{Step 4: $\tilde{D}_i$ is Already Demeaned}
%--------------------------------------------------------------------

We verify $E[\tilde{D}_i] = 0$ using the Law of Iterated Expectations:
\[
  E[\tilde{D}_i]
  = E\!\left[E[D_i - p(X_i) \mid X_i]\right]
  = E\!\left[E[D_i \mid X_i] - p(X_i)\right]
  = E\!\left[p(X_i) - p(X_i)\right]
  = 0.
\]
Since $E[\tilde{D}_i] = 0$, it follows immediately that
\[
  \tilde{D}_i - E[\tilde{D}_i] \;=\; \tilde{D}_i - 0 \;=\; \tilde{D}_i \;=\; D_i - p(X_i).
\]
The partialled-out variable is its own demeaned version.

%--------------------------------------------------------------------
\section{Summary}
%--------------------------------------------------------------------

\begin{proposition}
Let $X_i \in \{0,1\}$ be binary and $p(x) = E[D_i \mid X_i = x]$ be the
propensity score.  In the saturated OLS regression of $D_i$ on $(1, X_i)$:
\begin{enumerate}
  \item The constant and slope equal $p(0)$ and $p(1)-p(0)$ respectively.
  \item The fitted value equals $p(X_i)$ for every observation.
  \item The residual (partialled-out treatment) is $\tilde{D}_i = D_i - p(X_i)$.
  \item $E[\tilde{D}_i] = 0$, so $\tilde{D}_i$ is already mean-zero.
\end{enumerate}
\end{proposition}

\begin{remark}
The same argument extends to a discrete $X_i$ with $K$ cells: a fully saturated
regression includes a dummy for each cell, the fitted value in each cell is the
cell mean $p(x)$, and the residual is again $D_i - p(X_i)$.  This is the
foundation of the Frisch--Waugh--Lovell result used in Angrist (1998).
\end{remark}

\end{document}
